% checklist_pass.tex — 合规迷你论文 fixture（Wave6-A）
\documentclass{article}
\usepackage{booktabs}
\begin{document}

\section{摘要}
针对问题1，本文采用优化方法对系统进行建模，首先建立目标函数，接着引入约束条件，最后采用启发式算法求解，得到最优方案。针对问题2，本文在问题1的基础上扩展约束，采用灵敏度分析验证稳定性。实验表明本文方法性能优越，求解结果稳定可靠。
关键词：优化；灵敏度分析；启发式算法；约束建模

\section{引言}
\textbf{问题1}：在给定资源约束下，求解最小成本方案。
\textbf{问题2}：在问题1基础上，扩展多目标约束。

\section{总体分析}
本文围绕两个问题展开分析，流程图如图\ref{fig:flow}所示。
\begin{figure}[htbp]
  \centering
  \includegraphics[width=0.8\textwidth]{flow.pdf}
  \caption{总体流程图}
  \label{fig:flow}
\end{figure}

\section{模型假设}
为简化问题，本文做出以下假设：
假设1：系统处于稳态运行。
假设2：外部扰动服从正态分布。

\section{符号说明}
本文使用的主要符号如表\ref{tab:sym}所示。
\begin{table}[htbp]
  \centering
  \caption{主要符号说明}
  \label{tab:sym}
  \begin{tabular}{cl}
    \toprule
    符号 & 含义 \\
    \midrule
    $x$ & 决策变量 \\
    $c$ & 单位成本 \\
    \bottomrule
  \end{tabular}
\end{table}

\section{模型求解}
经过建模与求解，得到结果如表\ref{tab:res}与图\ref{fig:res}所示。
\begin{table}[htbp]
  \centering
  \caption{求解结果对比}
  \label{tab:res}
  \begin{tabular}{lc}
    \toprule
    方案 & 成本（元） \\
    \midrule
    方案一 & 1200 \\
    方案二 & 980 \\
    \bottomrule
  \end{tabular}
\end{table}
如图\ref{fig:res}所示，方案二在成本上优于方案一。
\begin{figure}[htbp]
  \centering
  \includegraphics[width=0.8\textwidth]{result.pdf}
  \caption{成本对比柱状图}
  \label{fig:res}
\end{figure}

\section{模型检验}
本文对优化模型进行灵敏度分析，结果稳定。

\section{模型的评价}
优点1：建模简洁，求解效率高。
优点2：灵敏度分析充分，结果稳健。
缺点1：未考虑动态扰动。

\section{模型的改进}
未来可引入动态扰动项，提升模型适应性。

\section{模型的推广}
本文方法可推广到同类资源调度问题。

\begin{thebibliography}{9}
\bibitem{a} 张三. 优化方法研究[J]. 系统工程理论与实践, 2023, 43(2): 1-10.
\bibitem{b} 李四. 灵敏度分析综述[M]. 北京: 科学出版社, 2022.
\end{thebibliography}

\end{document}
